\section{Introduction}
\label{sec:intro}

Automated recognition of individual student engagement from classroom
video would underpin every form of scalable learning analytics ---
from automated formative feedback to teacher dashboards that flag
at-risk learners in real time. The challenge: the standard benchmark
for this task --- \textbf{DAiSEE}~\cite{gupta2016daisee}, a
subject-disjoint dataset of 8{,}925 ten-second clips of 112 students
labelled with ordinal four-class engagement --- is widely cited but
poorly characterised in the frozen-feature regime that deployment
systems increasingly use. The standard linear-probe baseline on the
strongest publicly-available pre-trained encoder we tested
(SigLIP-L~\cite{zhai2023siglip}) achieves
$\kappa_q^{\text{argmax}} = 0.199$ --- enough to be useful, but no
prior paper has reported a careful recipe-tuned upper bound on this
regime.

\textbf{Our headline contribution} is \textbf{SETA} (Subject-aware
Engagement via Threshold-tuned Aggregation), the strongest
frozen-feature DAiSEE engagement recipe we are aware of:
\begin{enumerate}\itemsep0pt
  \item Encode the mid-clip frame ($t = 5$~s) with SigLIP-L
   (1024-d CLS).
  \item Train $K = 20$ bootstrap-resampled class-balanced logistic
   regression probes over five outer seeds (100 probes total),
   $C$-tuned per probe on validation $\kappa_q$, and average their
   predicted probabilities.
  \item Compute the expected-class score
   $\mathbb{E}[y] = \sum_{k} k \cdot p_k$ on validation and test.
  \item On validation, exhaustively search a grid of three ordinal
   thresholds $0 \le t_1 < t_2 < t_3 \le 3$ (step $0.04$) and select
   the triple maximising val $\kappa_q$; apply the same thresholds
   on test.
\end{enumerate}

SETA achieves \textbf{$\kappa_q = 0.247~[0.201, 0.294]$} (95\%
bootstrap CI, $n = 1{,}000$ resamples) on the full 1{,}784-clip
subject-disjoint DAiSEE test split --- a \textbf{$+0.048$ absolute
($+24\%$ relative)} improvement over the standard single-LR argmax
baseline, with the lower CI 0.201 statistically above the baseline
point estimate 0.199. The recipe runs end-to-end in under thirty
minutes on a single CPU; no encoder fine-tuning, no proprietary
inputs.

\textbf{Why does the ceiling sit here?} The natural follow-up
question is whether some readout-level or encoder-level trick
exceeds SETA. We investigate this systematically and find a
\emph{structural} answer.

\paragraph{The IDEP entanglement curve.} A class-balanced linear
classifier recovers subject identity from frozen SigLIP-L features
at \textbf{99.5\%} accuracy. We trace the engagement and identity
subspaces by ranking the identity classifier's weight rows by L2
norm and projecting features onto the orthogonal complement of the
top-$k$ identity directions. The result --- the
\textbf{identity-engagement entanglement (IDEP) curve}, our first
diagnostic contribution --- exhibits an asymmetric structure:
identity is encoded across dozens of redundant directions
(remaining recoverable at $> 95\%$ even after removing the top 20);
engagement variance is concentrated in only two-to-five
directions that overlap with the top of the identity basis.
Linear concept erasure~\cite{belrose2023leace} of identity
collapses engagement $\kappa_q$ to $0.04$ alongside it.

\paragraph{Within-subject ranking is essentially zero.} Our second
diagnostic shows that SETA's global $\kappa_q = 0.247$ hides a
striking behavioural pattern: within each test subject, the
Spearman rank correlation between predicted and true engagement
is \textbf{$\rho \approx 0.05$}. SETA --- and every other method we
tested --- captures \emph{which student} is typically more engaged
(subject-level priors) but cannot rank \emph{which of that
student's clips} is more engaged.

\paragraph{Cross-task control.} The same frozen encoders that
plateau at $\kappa_q \le 0.25$ on DAiSEE engagement reach
$\kappa = 0.55$ on FER2013 emotion classification --- bounding our
claim: the DAiSEE ceiling is engagement-task-specific, not
encoder-generic.

\paragraph{DREAM as constructive falsification.} The natural
remedy for identity dominance is per-subject personalization. We
introduce \textbf{DREAM}, a self-supervised neutral-anchor
personalization scheme: for each subject, select $K$
low-blendshape-activation frames (engagement-agnostic neutral
anchors, via MediaPipe FaceLandmarker~\cite{lugaresi2019mediapipe}),
compute the per-subject anchor embedding as the mean SigLIP-L
feature over those frames, and probe engagement on the residual.
DREAM is the gentlest linear identity intervention we can construct
--- yet it fails (best DREAM-cat $K{=}5$ bagged $\kappa_q = 0.218$,
below SETA), for an explainable reason that directly mirrors the
LEACE result and confirms the structural diagnosis: even gentle
linear identity removal disturbs the entangled engagement subspace.

\paragraph{Contributions.} We make five contributions:
\begin{enumerate}\itemsep0pt
  \item \textbf{SETA}: the strongest frozen-feature DAiSEE recipe in
   the literature, $\kappa_q = 0.247~[0.201, 0.294]$, $+24\%$
   relative over the argmax baseline, with bootstrap-significant
   improvement. Single-encoder, parameter-light, single-CPU
   reproducible.
  \item \textbf{IDEP entanglement curve}: a quantitative diagnostic
   that exposes the high-rank redundant structure of subject
   identity vs.\ the low-rank fragile structure of engagement in
   frozen features (Sec.~\ref{sec:idep}).
  \item \textbf{Within-subject ranking diagnostic}: a behavioural
   characterization showing that even at SETA's global $\kappa_q$,
   per-clip discrimination within a single subject is essentially
   random (Sec.~\ref{sec:within}).
  \item \textbf{DREAM}: a self-supervised neutral-anchor
   personalization scheme that constructively falsifies the
   ``remove identity, recover engagement'' hypothesis at the
   readout level (Sec.~\ref{sec:dream}).
  \item \textbf{Comprehensive method-family decomposition}: more
   than sixty method variants across eight encoder/data adaptation
   families (zero-shot, linear probe, ordinal, calibration,
   contrastive, ensemble, adapter, encoder fine-tune), all
   converging at $\kappa_q \in [0.20, 0.25]$ (Sec.~\ref{sec:results}).
\end{enumerate}

\paragraph{Reproducibility.} All extracted features, neutral
anchors, the full method-family grid, scripts, and bootstrap CIs
are released as supplementary material.
